\chapter{Insulin}

\section*{What Is This Test?}
Insulin is a 51-amino-acid polypeptide hormone produced by the beta cells of the pancreatic islets of Langerhans. It is synthesized as preproinsulin, cleaved to proinsulin (composed of an A chain, B chain, and connecting C-peptide), and then further cleaved in secretory granules to equimolar amounts of insulin and C-peptide, which are co-secreted into the portal circulation. Insulin is the primary anabolic hormone: it promotes glucose uptake in skeletal muscle and adipose tissue (via translocation of GLUT4 glucose transporters to the cell surface), stimulates glycogen synthesis in liver and muscle, promotes lipogenesis and inhibits lipolysis in adipose tissue, stimulates protein synthesis, and promotes cell growth and differentiation. Insulin secretion is stimulated primarily by glucose (via GLUT2-mediated uptake, glycolysis, and ATP-sensitive potassium channel closure leading to calcium influx and insulin granule exocytosis), with amplification by incretins (GLP-1 and GIP), amino acids, and vagal stimulation. Insulin is inhibited by hypoglycemia, somatostatin, amylin, and sympathetic nervous system activation (alpha-adrenergic). Insulin has a short plasma half-life of 4--6 minutes and undergoes approximately 50--60\% first-pass hepatic extraction. The insulin blood test measures the serum concentration of insulin, but its clinical utility as a standalone test is limited. Insulin is most useful when interpreted in the context of a simultaneously drawn glucose level, in the evaluation of hypoglycemia (Whipple triad: symptoms of hypoglycemia, low plasma glucose <55 mg/dL, and relief of symptoms with glucose administration), insulin resistance assessment, and research applications.

\section*{Alternative Names}
Serum Insulin; Plasma Insulin; Fasting Insulin; Immunoreactive Insulin; Insulin by Immunoassay; Insulin-Glucose Ratio; HOMA-IR (Homeostatic Model Assessment of Insulin Resistance)

\section*{Principle}
Insulin is measured using a two-site sandwich immunometric (chemiluminescent) immunoassay on automated analyzers. Two monoclonal antibodies against different epitopes on the insulin molecule are used: a capture antibody on a solid phase and a detection antibody labeled with a chemiluminescent tag (ruthenium on Roche ECLIA, acridinium ester on Abbott/Siemens). The assay must have minimal cross-reactivity with proinsulin and its conversion intermediates (des-31,32 and des-64,65 split proinsulins) because these are co-secreted and can accumulate in insulinoma and renal failure. Third-generation insulin assays have <0.1\% cross-reactivity with intact proinsulin. The functional sensitivity is typically <1--2 microIU/mL. The assay is calibrated against the WHO 1st International Reference Preparation for human insulin (IRP 66/304) \cite{tietz2023textbook}. Because hemolysis releases insulin-degrading enzyme from erythrocytes, hemolyzed specimens yield falsely low insulin values. Insulin measurement for hypoglycemia evaluation should be performed simultaneously with glucose and C-peptide, and with a specimen that is centrifuged and separated promptly to minimize in-vitro insulin degradation.

\section*{History}
Insulin was discovered in 1921 by Frederick Banting, Charles Best, and James Collip at the University of Toronto under the direction of John Macleod. Banting and Macleod were awarded the Nobel Prize in 1923. The first insulin injection was administered to 14-year-old Leonard Thompson on January 11, 1922 at Toronto General Hospital. The amino acid sequence of insulin was determined by Frederick Sanger in 1955 (Nobel Prize 1958). The first radioimmunoassay for insulin was developed by Rosalyn Yalow and Solomon Berson in 1959 at the Bronx VA Hospital (Nobel Prize 1977 for Yalow) \cite{yalow1960immunoassay}. The development of recombinant human insulin by Genentech in 1978 (Humulin, marketed 1982) eliminated reliance on animal insulin. Insulin analogues (lispro, 1996; glargine, 2000) were developed to improve pharmacokinetic profiles. The role of insulin measurement in insulinoma diagnosis was established in the 1970s--1980s. The HOMA-IR model for estimating insulin resistance was developed by Matthews and colleagues in 1985 \cite{matthews1985homeostasis}.

\section*{Why This Test Is Ordered}
\begin{itemize}
  \item Evaluation of fasting hypoglycemia (Whipple triad) to diagnose insulinoma --- inappropriately elevated insulin (>3 microIU/mL, often >5--10 microIU/mL) with concurrent hypoglycemia (glucose <55 mg/dL) and elevated C-peptide confirms endogenous hyperinsulinism. Insulin-to-glucose ratio >0.3 is suggestive of insulinoma \cite{service1995hypoglycemic}.
  \item Diagnosis of insulinoma --- a rare insulin-secreting pancreatic neuroendocrine tumor (incidence 1--4 per million per year). Fasting hypoglycemia with inappropriately normal or elevated insulin during a supervised 72-hour fast \cite{deherder2006welldifferentiated}.
  \item Investigation of postprandial (reactive) hypoglycemia --- mixed meal test evaluating glucose and insulin responses at 0, 30, 60, 120, 180, 240, and 300 minutes post-meal
  \item Assessment of insulin resistance in metabolic syndrome, PCOS, and NAFLD --- fasting insulin >15--20 microIU/mL with normal glucose suggests insulin resistance. HOMA-IR calculated as (fasting insulin x fasting glucose) / 405
  \item Detection of factitious hypoglycemia from surreptitious exogenous insulin or sulfonylurea ingestion --- elevated insulin (>50--100 microIU/mL can be seen) with suppressed C-peptide (<0.6 ng/mL) indicates exogenous insulin administration; elevated insulin and C-peptide indicates sulfonylurea-induced endogenous hyperinsulinism
  \item Evaluation of type 2 diabetes pathophysiology and monitoring of beta cell function over time
  \item Assessment of insulin requirements and beta cell reserve in type 1 and type 2 diabetes --- C-peptide is preferred for this indication because of hepatic extraction of insulin
  \item Research in metabolic syndrome, obesity, and cardiovascular risk
\end{itemize}

\section*{Primary vs Secondary Test}
Insulin is a secondary or specialized test. For hypoglycemia evaluation, the primary tests are plasma glucose, C-peptide, and proinsulin, with insulin being an essential co-test. For insulin resistance assessment, fasting insulin and HOMA-IR are secondary (surrogate) markers used in research and specialized clinical settings; fasting glucose and hemoglobin A1c are primary. Insulin measurement during a supervised 72-hour fast is the gold standard for insulinoma diagnosis but is a specialized procedure performed at tertiary centers \cite{cryer2009evaluation}.

\section*{How This Test Is Performed}
A healthcare professional draws a fasting blood sample from a vein into a serum separator tube (gold-top) or plain red-top tube (~3--5 mL). For concurrent glucose measurement, a fluoride oxalate (gray-top) tube is also drawn because glycolysis in uncentrifuged blood lowers glucose by 5--7\% per hour at room temperature. For hypoglycemia evaluation, the sample must be drawn at the time of spontaneous hypoglycemia (glucose <55 mg/dL) and sent STAT for glucose, insulin, C-peptide, proinsulin, beta-hydroxybutyrate, and sulfonylurea/meglitinide screen. For insulin resistance assessment, a morning fasting sample (8--10 hours) is drawn. In the 72-hour fast for insulinoma diagnosis, the patient is hospitalized, allowed only water and non-caloric beverages, and glucose is measured every 4--6 hours; when glucose drops below 60 mg/dL, frequency increases to every 1--2 hours; when glucose drops below 45 mg/dL with symptoms, the critical blood sample is drawn and the fast is terminated with IV dextrose.

\section*{What Preparation Is Needed}
Fasting for 8--12 hours is required for fasting insulin assessment. For hypoglycemia evaluation, the sample must be drawn during spontaneous hypoglycemia when symptoms are present. Strenuous exercise should be avoided for 24 hours. High-dose biotin (>5 mg/day) must be held 48 hours before testing. Medications that affect insulin secretion should be documented: sulfonylureas, meglitinides (increase); diazoxide, octreotide, calcium channel blockers (decrease); exogenous insulin administration (markedly elevates measured insulin, especially with insulin analogue cross-reactivity --- some insulin analogues cross-react weakly with insulin immunoassays). For the 72-hour fast, all medications affecting glucose metabolism should be held.

\section*{Contraindications and Precautions}
\begin{itemize}
  \item The most critical limitation is that insulin measured from a peripheral vein reflects only 40--50\% of portal insulin secretion because of first-pass hepatic extraction (approximately 50--60\% of portal insulin is removed by the liver). C-peptide, which is co-secreted in equimolar amounts with insulin but does not undergo hepatic extraction, is a more accurate marker of endogenous insulin secretion. For this reason, C-peptide is preferred for assessing beta cell function. Insulin immunoassays cross-react to varying degrees with proinsulin (5--40\% in some assays) and insulin analogues (lispro, aspart, glargine, detemir: 0--80\% depending on the analogue and assay). This is important: if testing for surreptitious insulin administration, the specific insulin analogue being used may not cross-react well, leading to falsely low total insulin despite clinical hypoglycemia. Hemolysis falsely lowers insulin because erythrocytes release insulin-degrading enzyme (IDE). Insulin is stable at room temperature for 8 hours. In insulin resistance, fasting insulin >15--20 microIU/mL with normal glucose is suggestive, but insulin levels must be interpreted with glucose simultaneously. Anti-insulin antibodies in patients receiving exogenous insulin (especially older animal insulins) can interfere with immunoassays, causing falsely elevated or lowered results.
\end{itemize}
\section*{Risks}
Minimal risk from venipuncture. For the 72-hour fast: prolonged fasting, discomfort, and risk of severe symptomatic hypoglycemia requiring IV dextrose. The fast MUST be conducted under medical supervision with IV access.

\section*{What Happens After the Test}
Results available within 1--4 hours on automated analyzers (insulin) or 3--7 days for proinsulin and sulfonylurea screen. For hypoglycemia evaluation: Whipple triad confirmed (symptoms, glucose <55 mg/dL, relief with glucose) plus elevated insulin (>3 microIU/mL), elevated C-peptide (>0.6 ng/mL), elevated proinsulin, and suppressed beta-hydroxybutyrate (<2.7 mmol/L) confirms endogenous hyperinsulinism. Next step: localization with abdominal CT/MRI, endoscopic ultrasound, and possibly selective arterial calcium stimulation with hepatic venous sampling. For insulin resistance: fasting insulin >15--20 microIU/mL and HOMA-IR >2.5--3.0 suggests insulin resistance. For surreptitious insulin use: elevated insulin with suppressed C-peptide confirms exogenous insulin administration.

\section*{Understanding Results}

\subsection*{Below Normal}
\begin{itemize}
  \item Type 1 diabetes mellitus: absolute insulin deficiency due to autoimmune destruction of beta cells. Fasting insulin typically <3 microIU/mL (or undetectable) with hyperglycemia. C-peptide is also low/undetectable.
  \item Advanced type 2 diabetes: beta cell failure after years of insulin resistance. Low or normal fasting insulin despite hyperglycemia indicates beta cell exhaustion.
  \item Fasting (prolonged), starvation, or ketogenic diet: insulin is appropriately suppressed to <3--5 microIU/mL to permit lipolysis, ketogenesis, and hepatic glucose production.
  \item Diazoxide or octreotide therapy: pharmacologic suppression of insulin secretion. Diazoxide opens beta cell potassium channels, preventing depolarization; octreotide (somatostatin analogue) directly suppresses secretion.
  \item Post-pancreatectomy: removal of beta cells eliminates insulin production. Patients require lifelong exogenous insulin.
\end{itemize}

\subsection*{Above Normal}
\begin{itemize}
  \item Insulinoma: inappropriately elevated fasting insulin (>3 microIU/mL, often >5--10 microIU/mL) with concurrent hypoglycemia (glucose <55 mg/dL). The insulin-to-glucose ratio is typically >0.3 (insulin in microIU/mL divided by glucose in mg/dL). During the 72-hour fast, insulin fails to suppress appropriately: normal response to hypoglycemia is insulin <3 microIU/mL. Elevated insulin with elevated C-peptide and proinsulin confirms endogenous hyperinsulinism (insulinoma). Negative sulfonylurea/meglitinide screen rules out factitious sulfonylurea-induced hypoglycemia.
  \item Insulin resistance (metabolic syndrome, type 2 diabetes, obesity, PCOS): compensatory hyperinsulinemia. Fasting insulin typically 15--50 microIU/mL with normal or elevated fasting glucose. Insulin resistance underlies the metabolic syndrome cluster: abdominal obesity, hypertension, dyslipidemia (high triglycerides, low HDL), and impaired glucose regulation.
  \item Sulfonylurea or meglitinide-induced hyperinsulinism: these drugs close beta cell potassium channels, stimulating insulin secretion. In surreptitious use, presents identically to insulinoma with elevated insulin, C-peptide, and proinsulin. Positive sulfonylurea screen on the critical blood sample confirms the diagnosis. Requires toxicology or specialized laboratory testing.
  \item Exogenous insulin administration (therapeutic or factitious): total insulin markedly elevated (often >50--100 microIU/mL, may exceed 1000 microIU/mL depending on dose) with C-peptide suppressed (<0.6 ng/mL). Insulin:C-peptide molar ratio >1.0 indicates exogenous insulin. This is the key distinction from insulinoma where both insulin and C-peptide are elevated.
  \item Insulin autoimmune syndrome (Hirata disease): autoantibodies against endogenous insulin cause postprandial hyperglycemia followed by late hypoglycemia (insulin-antibody complexes dissociate hours after meals, releasing free insulin). Markedly elevated total insulin (often >1000 microIU/mL) with normal C-peptide. Associated with HLA-DR4 and exposure to sulfhydryl-containing drugs (methimazole, captopril, penicillamine). Anti-insulin antibodies confirm.
  \item Congenital hyperinsulinism (nesidioblastosis): neonatal or infantile persistent hypoglycemia from diffuse or focal beta cell hyperplasia due to mutations in genes regulating insulin secretion (ABCC8, KCNJ11 encoding SUR1 and Kir6.2 subunits of the ATP-sensitive potassium channel --- most common). Inappropriately elevated insulin with hypoglycemia requiring glucose infusion, diazoxide, or octreotide.
  \item Post-gastric bypass hypoglycemia (dumping syndrome): rapid glucose delivery to the small intestine stimulates excessive GLP-1 and insulin secretion, causing late postprandial hypoglycemia 1--4 hours after meals. Insulin is inappropriately elevated at the time of hypoglycemia. Treatment: dietary modification (low-glycemic index, smaller meals), acarbose, octreotide.
\end{itemize}

\section*{Normal Results}
\begin{itemize}
  \item Fasting insulin (euglycemic, insulin-sensitive): 2--12 microIU/mL (14--86 pmol/L)
  \item Fasting insulin in insulin resistance: 15--30+ microIU/mL
  \item Insulin-to-glucose ratio: <0.3 (normal fasting); >0.3 suggests insulinoma (in setting of hypoglycemia)
  \item HOMA-IR = (fasting insulin x fasting glucose) / 405: <1.0 (insulin sensitive), 1.0--2.5 (normal), >2.5 (insulin resistance), >3.0 (significant insulin resistance)
  \item Insulin in insulinoma (during hypoglycemia): >3 microIU/mL (inappropriate for glucose <55 mg/dL)
  \item Insulin in exogenous insulin administration: >50--100 microIU/mL (often much higher)
  \item Insulin:C-peptide molar ratio: <1.0 (endogenous); >1.0 (exogenous insulin)
  \item Note: Reference ranges are assay- and laboratory-specific. Insulin must always be interpreted with simultaneous glucose and preferably C-peptide. Hemolyzed specimens are rejected because of insulin degradation.
\end{itemize}

\section*{Follow-up Tests}
\begin{itemize}
  \item Simultaneous plasma glucose (fluoride oxalate tube): Must be drawn at the same time as insulin. Insulin:glucose ratio and HOMA-IR require both.
  \item C-peptide: The most important co-test. Distinguishes endogenous hyperinsulinism (elevated C-peptide) from exogenous insulin administration (suppressed C-peptide). Also the preferred marker for beta cell function because it avoids hepatic extraction.
  \item Proinsulin: Elevated in insulinoma (>5--20 pmol/L), reflecting incomplete processing of proinsulin to insulin by tumor cells. Normal proinsulin is <20\% of total insulin immunoreactivity. Proinsulin >30\% is suggestive of insulinoma.
  \item Beta-hydroxybutyrate (BHB): Suppressed in insulin-mediated hypoglycemia (<2.7 mmol/L) because insulin inhibits lipolysis and ketogenesis. Elevated in non-insulin-mediated hypoglycemia (starvation, adrenal insufficiency, alcohol).
  \item Sulfonylurea/meglitinide screen (toxicology): Essential in all cases of confirmed hyperinsulinemic hypoglycemia to exclude surreptitious sulfonylurea or meglitinide use before proceeding to insulinoma localization. Plasma or urine toxicology screen.
  \item Anti-insulin antibodies: Diagnose insulin autoimmune syndrome (Hirata disease) in patients with markedly elevated insulin without known exogenous insulin exposure.
  \item Abdominal CT or MRI (pancreatic protocol, triple-phase): To localize insulinoma. Insulinomas are typically small (<2 cm, 80--90\%), hypervascular, well-defined pancreatic lesions. CT may miss up to 30--40\% of insulinomas.
  \item Endoscopic ultrasound (EUS): The most sensitive preoperative localization technique for pancreatic insulinomas. Sensitivity 85--95\%. Allows FNA for cytology.
  \item Selective arterial calcium stimulation with hepatic venous sampling: The gold standard for regionalizing insulinoma when non-invasive imaging is negative. Calcium gluconate is injected into the splenic, gastroduodenal, and superior mesenteric arteries sequentially; hepatic venous insulin is measured. A >2--3 fold rise in insulin after calcium injection into a specific artery localizes the insulinoma to that arterial territory.
  \item Ga-68 DOTATATE or Exendin-4 PET/CT: Somatostatin receptor imaging (if DOTATATE) or GLP-1 receptor imaging (Exendin-4, more specific for insulinoma that expresses GLP-1 receptors) for tumor localization when conventional imaging is negative.
\end{itemize}

\section*{References}
\bibliographystyle{unsrtnat}
\bibliography{references}

\clearpage
\section*{Regulatory Status and Authorization Date}
This chapter describes a test category, not a specific commercial product. FDA, CDSCO, EU IVDR, and MHRA approval or registration dates therefore cannot be assigned without the assay manufacturer, product name, instrument, and intended use. For laboratory-developed tests, an individual agency approval date may not exist. See the Preface for the interpretation of regulatory status.

\clearpage
